\documentclass[11pt]{article}
\usepackage[margin=1in]{geometry}
\usepackage{amsmath,amsthm,mathtools}
\usepackage{amssymb}

\newtheorem{theorem}{Theorem}[section]
\newtheorem{lemma}[theorem]{Lemma}
\newtheorem{proposition}[theorem]{Proposition}
\newtheorem{corollary}[theorem]{Corollary}
\theoremstyle{definition}
\newtheorem{definition}[theorem]{Definition}
\newtheorem{remark}[theorem]{Remark}

\newcommand{\R}{\mathbb{R}}
\newcommand{\Z}{\mathbb{Z}}
\newcommand{\T}{\mathbb{T}}
\newcommand{\cC}{\mathcal{C}}
\DeclareMathOperator{\re}{Re}
\DeclareMathOperator{\im}{Im}

\title{Toward $g^{(1)}\geq c > 0$: Proof Architecture\\
{\large Working notes --- not for circulation}}

\author{R.\ Fontes \& Claude (Anthropic)\\
Quijotic Research}

\date{March 2026}

\begin{document}
\maketitle

%======================================================================
\section{The target}
%======================================================================

\subsection{Setup}

Let $u$ be a smooth divergence-free field on $\T^3$
with purely imaginary Fourier structure:
$\hat u(k) = iV(k)$ where $V(k)\in\R^3$,
$k\cdot V(k)=0$, and $V(-k)=-V(k)$.
Partition $S^2$ into $J$ sectors $\{\Omega_j\}$.
Define:
\begin{align}
E_j^{(1)} &= \sum_{\hat k\in\Omega_j}|k|^2|V(k)|^2,
\qquad E^{(1)} = \sum_j E_j^{(1)}, \qquad
p_j = E_j^{(1)}/E^{(1)}, \\
\cC^{(1)} &= \sum_j p_j^2, \\
T^{(1)}(k) &= 2|k|^2\sum_{p+q=k}
(q\cdot V(p))(V(k)\cdot V(q)),
\label{eq:T1}\\
\Sigma^{(1)} &= \frac{2}{E^{(1)}}
\sum_j p_j\Bigl(\sum_{\hat k\in\Omega_j}T^{(1)}(k)
- p_j\sum_k T^{(1)}(k)\Bigr).
\label{eq:Sigma1}
\end{align}
Note: for the purely imaginary class,
$T^{(1)}(k)$ is real (the $\im$ from the general
formula collapses because $z(k;p,q) =
i(q\cdot V(p))(V(k)\cdot V(q))$ is purely
imaginary, so $\im[z]$ is the real product).

\subsection{The claim}

\begin{equation}\label{eq:target}
\boxed{g^{(1)}(u) :=
\frac{-\Sigma^{(1)}(u)}{\cC^{(1)}(u)-1/J}
\geq c > 0}
\end{equation}
for all divergence-free $V$ with
$\cC^{(1)}>1/J$ and some universal $c$
depending only on $J$ (and possibly on the
support of $V$).

Numerics give $g^{(1)}\approx 0.47\pm 0.04$
for TG-generated fields with
$\cC^{(1)}\in[0.13, 0.50]$ at $J=32$.


%======================================================================
\section{Proof architecture}
%======================================================================

The proof proceeds in three stages of
increasing generality.

\subsection{Stage 1: Two-sector model (this document)}

Restrict to fields where enstrophy lives
in exactly two sectors: one dominant
($\Omega_*$ with fraction $p_*$) and one
minority ($\Omega_\dagger$ with fraction
$1-p_*$).  In this setting, $\Sigma^{(1)}$
reduces to a single sum over cross-sector
triads, and the divergence-free constraint
$k\cdot V(k)=0$ creates a computable
geometric bias.  The goal is to prove
$g^{(1)}>0$ with an explicit lower bound.

\subsection{Stage 2: Dominant-sector
decomposition}

For general $J$-sector fields, decompose
$\Sigma^{(1)}$ into:
\begin{equation}
\Sigma^{(1)} = \Sigma^{(1)}_{\rm cross}
+ \Sigma^{(1)}_{\rm within}.
\end{equation}
The numerics (N7) show
$|\Sigma_{\rm cross}|/|\Sigma_{\rm within}|
\approx 830$.  Prove:
\begin{enumerate}
\item[(a)] $\Sigma^{(1)}_{\rm cross}<0$
  with $|\Sigma^{(1)}_{\rm cross}|
  \geq c_1(\cC^{(1)}-1/J)$.
\item[(b)] $|\Sigma^{(1)}_{\rm within}|
  \leq c_2(\cC^{(1)}-1/J)$ with $c_2\ll c_1$.
\end{enumerate}
This reduces the problem to bounding the
cross-sector transfer, where the geometric
bias from the divergence-free constraint is
strongest.

\subsection{Stage 3: Trilinear combinatorics}

Express $\Sigma^{(1)}_{\rm cross}$ as a sum
over lattice triads weighted by $|k|^2$ and
constrained by $k\cdot V(k)=0$.  Bound the
sum using the angular separation between
sectors (which controls the geometric bias)
and the Cauchy--Schwarz inequality on the
triad factors.


%======================================================================
\section{Stage 1: The two-sector lemma}
%======================================================================

\subsection{The reduced problem}

Suppose $V$ is supported on wavevectors
whose directions $\hat k$ lie in exactly two
sectors $\Omega_*$ and $\Omega_\dagger$, with
enstrophy fractions $p_*>1/2$ and
$p_\dagger = 1-p_*$.  Then
$\cC^{(1)} = p_*^2 + p_\dagger^2$ and
$\cC^{(1)}-1/2 = 2(p_*-1/2)^2\geq 0$
(using $J=2$ for the two-sector model).

The spreading functional reduces to
\begin{equation}\label{eq:Sigma-2sector}
\Sigma^{(1)} = \frac{2}{E^{(1)}}
\bigl[p_*(T_*-p_*T_{\rm tot})
+ p_\dagger(T_\dagger - p_\dagger T_{\rm tot})\bigr],
\end{equation}
where $T_j = \sum_{\hat k\in\Omega_j}T^{(1)}(k)$.
Since $p_*+p_\dagger=1$ and
$T_*+T_\dagger = T_{\rm tot}$, this simplifies to
\begin{equation}\label{eq:Sigma-simple}
\Sigma^{(1)} = \frac{2}{E^{(1)}}
(p_*-p_\dagger)(T_\dagger - p_\dagger T_{\rm tot})
= \frac{2(2p_*-1)}{E^{(1)}}
(T_\dagger - p_\dagger T_{\rm tot}).
\end{equation}
Since $p_*>1/2$, we have $2p_*-1>0$.
So $\Sigma^{(1)}<0$ iff
$T_\dagger < p_\dagger T_{\rm tot}$:
the minority sector receives less than its
proportional share of the total transfer.

Equivalently, defining the excess transfer
to the dominant sector:
\begin{equation}
\Delta := T_* - p_* T_{\rm tot}
= -(T_\dagger - p_\dagger T_{\rm tot}),
\end{equation}
we need $\Delta > 0$ when $p_* > 1/2$.
That is: the dominant sector receives
\emph{more} than its share of the total
nonlinear enstrophy transfer---but in the
``wrong direction'' (depletion).

Wait---let us be more careful.
$T^{(1)}(k) = 2|k|^2\re[\hat u^*\cdot\hat B]$
measures the rate of enstrophy input to mode $k$.
$T^{(1)}(k)>0$ means mode $k$ gains enstrophy;
$T^{(1)}(k)<0$ means it loses.  The spreading
condition $\Sigma^{(1)}<0$ means sectors with
above-average enstrophy fraction ($p_j > 1/J$)
tend to have below-average transfer.

In the two-sector case: $\Sigma^{(1)}<0$ iff
the high-$p$ sector $\Omega_*$ has
$T_*/E_*^{(1)} < T_{\rm tot}/E^{(1)}$,
i.e., the per-unit-enstrophy transfer rate
is lower in the concentrated sector.


\subsection{Triad decomposition}

The total transfer to the dominant sector is
\begin{equation}
T_* = \sum_{\hat k\in\Omega_*}
2|k|^2\sum_{p+q=k}
(q\cdot V(p))(V(k)\cdot V(q)).
\end{equation}
Decompose by the sector membership of $p$:
\begin{align}
T_*^{(\rm ww)} &= \sum_{\hat k\in\Omega_*}
2|k|^2\sum_{\substack{p+q=k\\
\hat p\in\Omega_*,\;\hat q\in\Omega_*}}
(q\cdot V(p))(V(k)\cdot V(q)),
\label{eq:Tww}\\
T_*^{(\rm cw)} &= \sum_{\hat k\in\Omega_*}
2|k|^2\sum_{\substack{p+q=k\\
\hat p\in\Omega_\dagger}}
(q\cdot V(p))(V(k)\cdot V(q)),
\label{eq:Tcw}
\end{align}
so $T_* = T_*^{(\rm ww)}+T_*^{(\rm cw)}$
(``within--within'' and ``cross--within'').

The key claim: the cross-sector contribution
$T_*^{(\rm cw)}$ is systematically negative
(depletion) due to the geometric bias from
$k\cdot V(k)=0$.


\subsection{The geometric bias}

\begin{lemma}[Divergence-free orthogonality
bias]\label{lem:geom}
Let $k,p\in\Z^3\setminus\{0\}$ with
$\hat k\in\Omega_*$ and $\hat p\in\Omega_\dagger$,
and let $q = k-p$.  Suppose $V(k)\perp k$ and
$V(p)\perp p$.  Let $\theta$ be the angle between
the sectors $\Omega_*$ and $\Omega_\dagger$
(i.e., between $\hat k$ and $\hat p$).

Then the triad factor satisfies:
\begin{equation}
(q\cdot V(p))(V(k)\cdot V(q))
= |V(p)||V(q)||q|
\cos\angle(q,V(p))\cos\angle(V(k),V(q)),
\end{equation}
and the divergence-free constraints create
a systematic sign bias:
\begin{enumerate}
\item $V(p)\perp p$ means $V(p)$ lies in the
  plane perpendicular to $\hat p$.  Since
  $q = k-p$ and $\hat k\not\parallel\hat p$,
  the projection $q\cdot V(p)$ depends on the
  component of $q$ in the plane $\perp p$.
\item Similarly, $V(k)\perp k$ constrains
  the second factor $V(k)\cdot V(q)$.
\end{enumerate}
\end{lemma}

The geometric content is clearest in the
\textbf{planar reduction}: restrict to
wavevectors $k = (k_1, k_2, 0)$ lying in
a plane, with $V(k) = |V|(- k_2, k_1, 0)/|k|$
(the unique unit divergence-free direction
in the plane, up to sign).


\subsection{Planar two-sector computation}

Let $\Omega_*$ contain wavevectors near
$\hat e_1 = (1,0,0)$ and $\Omega_\dagger$
contain wavevectors near $\hat e_2 = (0,1,0)$
(orthogonal sectors, $\theta=\pi/2$).

Take a minimal configuration:
$k = (n,0,0)\in\Omega_*$ and
$p = (0,m,0)\in\Omega_\dagger$ with $n,m>0$.
Then $q = k-p = (n,-m,0)$.

Divergence-free:
$V(k) = a_k(0,1,0)$ and
$V(p) = a_p(1,0,0)$ for some $a_k, a_p\in\R$.

Compute the triad factors:
\begin{align}
q\cdot V(p) &= (n,-m,0)\cdot a_p(1,0,0) = na_p,\\
V(k)\cdot V(q) &= a_k(0,1,0)\cdot V(q).
\end{align}
Now $V(q)$ for a mode at $q=(n,-m,0)$:
if $q$ is also a mode in $V$, then
$V(q)\perp q$, so
$V(q) = a_q(m,n,0)/\sqrt{n^2+m^2}$.
Then:
\begin{equation}
V(k)\cdot V(q) = a_k\cdot\frac{a_q\,n}
{\sqrt{n^2+m^2}}.
\end{equation}
The triad contribution is:
\begin{equation}\label{eq:planar-triad}
(q\cdot V(p))(V(k)\cdot V(q))
= \frac{n^2\,a_k\,a_p\,a_q}{\sqrt{n^2+m^2}}.
\end{equation}
This has a definite sign determined by
$\sgn(a_k\,a_p\,a_q)$.

The transfer to mode $k\in\Omega_*$ from
this cross-sector triad is:
\begin{equation}
T^{(1)}_{\rm cross}(k) = 2n^2\cdot
\frac{n^2\,a_k\,a_p\,a_q}{\sqrt{n^2+m^2}}.
\end{equation}

\begin{remark}[Sign from the TG structure]
For the TG vortex: $V(k)=(-k_1/8, k_2/8, 0)$,
so at $k=(1,1,1)$: $a_k$ involves mixed signs
from $k_1, k_2$.  The exact computation
(Theorem~4.1 of the paper) shows these signs
conspire to give $T_{\rm cross}<0$ at the
dominant modes.  The question is whether this
conspiracy is generic or specific to TG.
\end{remark}


\subsection{The sign mechanism}

The key observation: in the cross-sector
triad $(k,p,q)$ with $k\in\Omega_*$ and
$p\in\Omega_\dagger$, the three velocities
$V(k), V(p), V(q)$ are constrained by:
\begin{itemize}
\item $V(k)\perp k$ (div-free at $k$)
\item $V(p)\perp p$ (div-free at $p$)
\item $V(q)\perp q$ (div-free at $q$)
\item $q = k-p$ (triad constraint)
\end{itemize}
These four constraints, for well-separated
$\hat k$ and $\hat p$, create a
\textbf{geometric rigidity}: the product
$(q\cdot V(p))(V(k)\cdot V(q))$ has a
preferred sign that depends on the relative
orientation of the sectors.

\begin{proposition}[Sign lemma for
orthogonal sectors]\label{prop:sign}
Let $\Omega_*$ and $\Omega_\dagger$ be
orthogonal sectors (angular separation
$\theta = \pi/2$).  For any divergence-free
$V$ supported on modes in $\Omega_*\cup
\Omega_\dagger$ with uniform polarization
within each sector (i.e., $V(k)/|V(k)|$
constant on $\Omega_*$ and constant on
$\Omega_\dagger$), the cross-sector
transfer to $\Omega_*$ satisfies:
\begin{equation}
T_*^{(\rm cw)} = -C\sum_{\hat k\in\Omega_*}
|k|^2|V(k)|\sum_{\substack{p:\hat p\in
\Omega_\dagger\\k-p\neq 0}}
\frac{|k\cdot\hat e_p|^2\,|V(p)|\,|V(k-p)|}
{|k-p|}
\end{equation}
where $C>0$ is a geometric constant
depending on $\theta$.  In particular,
$T_*^{(\rm cw)}\leq 0$: cross-sector
transfer depletes the dominant sector.
\end{proposition}

\begin{remark}[The uniform polarization
hypothesis]
The uniform-polarization assumption is
strong: it says all modes in a sector share
the same velocity direction (up to amplitude).
This holds exactly for TG (where the 8 modes
have $V(k) = (-k_1/8, k_2/8, 0)$ and modes
in each sector are related by symmetry) and
approximately for concentrated fields (the
``narrow-sector hypothesis'' of Section~6.1
in the paper).  Removing it is the main
challenge for Stage~2.
\end{remark}


%======================================================================
\section{The within-sector bound}
%======================================================================

The within-sector transfer $T_*^{(\rm ww)}$
involves triads where all three legs are in
$\Omega_*$.  For such triads, $k$, $p$, and
$q=k-p$ all point in roughly the same
direction, so the divergence-free constraint
$V\perp k$ forces all three velocities to be
roughly parallel.  The product
$(q\cdot V(p))(V(k)\cdot V(q))$ then involves
$q\cdot V(p)$ where $q\approx k-p$ is
``nearly perpendicular'' to $V(p)$ (because
$V(p)\perp p$ and $q\approx k-p$ is roughly
parallel to $p$ within the same narrow sector).

This means within-sector triads have
$|q\cdot V(p)| \sim |V(p)|\cdot|q|\cdot
\sin(\text{opening angle})$,
which is suppressed by the sector opening
angle.  In contrast, cross-sector triads have
$|q\cdot V(p)| \sim |V(p)|\cdot|q|\cdot
\cos(\text{sector separation})$,
which is $O(1)$.

\begin{lemma}[Within-sector suppression]
\label{lem:within}
For a sector $\Omega_*$ with half-opening
angle $\alpha$:
\begin{equation}
|T_*^{(\rm ww)}| \leq C\alpha\cdot
\Bigl(\sum_{\hat k\in\Omega_*}
|k|^2|V(k)|^2\Bigr)^{3/2},
\end{equation}
while
\begin{equation}
|T_*^{(\rm cw)}| \geq c\cos\theta\cdot
\Bigl(\sum_{\hat k\in\Omega_*}
|k|^2|V(k)|^2\Bigr)
\Bigl(\sum_{\hat p\in\Omega_\dagger}
|p|^2|V(p)|^2\Bigr)^{1/2}.
\end{equation}
For $J$ equal-area sectors, $\alpha\sim
J^{-1/2}$, giving
$|T^{(\rm ww)}|/|T^{(\rm cw)}| = O(J^{-1/2})$.
This explains the numerical observation
$|\Sigma_{\rm cross}|/|\Sigma_{\rm within}|
\approx 830$ at $J=32$ (where
$J^{1/2}\approx 5.7$; the large ratio
reflects additional cancellation within
sectors).
\end{lemma}


%======================================================================
\section{What remains to prove}
%======================================================================

\subsection{For the two-sector lemma (Stage 1)}

\begin{enumerate}
\item \textbf{Remove the uniform polarization
  hypothesis.}  For modes in the same sector,
  $V(k)/|V(k)|$ varies as $k$ ranges over
  $\Omega_*$.  The variation is controlled by
  the sector opening angle: if
  $\hat k, \hat k'\in\Omega_*$, then
  $|V(k)/|V(k)| - V(k')/|V(k')||
  \leq C\alpha$.  This suggests a perturbation
  argument from the uniform case.

\item \textbf{Make the geometric constant
  explicit.}  The constant $C$ in
  Proposition~\ref{prop:sign} depends on the
  sector geometry.  For Fibonacci spiral
  partitions, the angular separation between
  adjacent sectors is $\sim J^{-1/2}$, which
  gives a lower bound on $\cos\theta$.

\item \textbf{Sum over all triad pairs.}
  The cross-sector depletion must be summed
  over all $(k,p)$ pairs, accounting for the
  $|k|^2$ weight and the triad constraint
  $k-p\in\Z^3\setminus\{0\}$.  This is a
  lattice convolution estimate.
\end{enumerate}

\subsection{For the general case (Stages 2--3)}

\begin{enumerate}
\item \textbf{Multi-sector decomposition.}
  With $J$ sectors, the dominant-sector
  analysis must account for transfer between
  all pairs $(\Omega_i, \Omega_j)$.  The key
  is that the sign bias from the divergence-free
  constraint applies to \emph{every}
  cross-sector pair, not just the dominant one.

\item \textbf{The polarization distribution.}
  In general, $V(k)/|V(k)|$ is not constant
  within a sector.  The averaging over
  polarizations within a sector produces
  additional cancellation that must be
  controlled.  The Cauchy--Schwarz inequality
  on the energy space gives
  $|\sum_k f(k)V(k)|^2 \leq
  (\sum_k|f(k)|^2)(\sum_k|V(k)|^2)$,
  which bounds the polarization fluctuations
  in terms of the enstrophy.

\item \textbf{Connection to the cascade
  structure.}  The numerics show $g^{(1)}$
  depends on the cascade phase structure
  (N19), not just on $\cC^{(1)}$.  This means
  the lower bound $c$ may depend on a
  measurable property of the field (e.g.,
  the fraction of enstrophy in the imaginary
  part of $\hat u$, which is $100\%$ for the
  purely imaginary class).  Characterizing
  which fields satisfy $g^{(1)}>0$ is
  ultimately the question of identifying the
  ``basin of attraction'' of the spreading
  mechanism.
\end{enumerate}


%======================================================================
\section{Immediate next computation}
%======================================================================

The most productive next step is a
\textbf{symbolic verification} of
Proposition~\ref{prop:sign} for a concrete
two-sector configuration on a small grid
(say $N=8$ or $N=16$), using exact rational
arithmetic as in the TG computation.

Specifically:
\begin{enumerate}
\item Choose two orthogonal sectors (e.g.,
  modes near $(1,0,0)$ vs.\ modes near
  $(0,1,0)$).
\item Assign divergence-free $V(k)$ with
  uniform polarization and arbitrary
  amplitudes in each sector.
\item Enumerate all cross-sector triads
  $(k,p,q)$ with $k\in\Omega_*$,
  $p\in\Omega_\dagger$.
\item Compute $T_*^{(\rm cw)}$ symbolically
  and verify the sign.
\item Repeat for non-uniform polarization
  (random $V(k)\perp k$) and check whether
  the sign persists.
\end{enumerate}

If the sign is robust under polarization
variation, that provides strong evidence for
Stage~2 and suggests the perturbation
argument will work.  If the sign fails for
some polarization configurations, that
identifies exactly which geometric condition
is needed.


\end{document}
